\documentclass[../main]{subfiles}
\usetikzlibrary{shapes,backgrounds}
\begin{document}
\chapter{Tarjetas Kanban y pizarras}
\img{images/02/kanban.jpg}{Tablero Kanban.}

\info{
Kanban}{
Se trata de un método visual de gestión de proyectos que permite a los equipos visualizar sus flujos de trabajo y la carga de trabajo. En un tablero Kanban, el trabajo se muestra en un proyecto en forma de tablero organizado por columnas. Tradicionalmente, cada columna representa una etapa del trabajo. 
}

\section{Introducción}

En el desarrollo de proyectos, es común utilizar herramientas que permitan organizar las tareas de manera efectiva. Una de estas herramientas es la pizarra o tarjeta Kanban, que permite visualizar el estado de las tareas en un proyecto y mejorar la gestión del flujo de trabajo.

\section{¿Qué es una Pizarra o Tarjeta Kanban?}

La pizarra o tarjeta Kanban es una herramienta visual para la gestión de proyectos que se originó en la industria manufacturera japonesa. Consiste en un panel o pizarra en la que se colocan tarjetas o post-its que representan las tareas que se deben realizar en un proyecto.

\section{Ventajas de las Pizarras o Tarjetas Kanban}

Entre las ventajas que ofrecen las pizarras o tarjetas Kanban, se encuentran:

\begin{itemize}
    \item Visualización del estado de las tareas en tiempo real.
    \item Identificación de cuellos de botella y áreas de mejora en el flujo de trabajo.
    \item Facilidad para priorizar tareas y asignar responsabilidades.
    \item Mejora en la comunicación y colaboración del equipo.
\end{itemize}

\section{Ejemplo de Uso de una Pizarra o Tarjeta Kanban}

A continuación, se muestra un ejemplo de uso de una pizarra o tarjeta Kanban para la gestión de un proyecto de desarrollo de software:


En este ejemplo, las tarjetas se dividen en cuatro columnas: Backlog, To Do, Doing y Done. Cada tarjeta representa una tarea que debe ser realizada


\actividad{ Responder, investigar}

\begin{enumerate}
    \item ¿Qué es Kanban y cómo funciona?
\item ¿Cuáles son los beneficios de utilizar Kanban en los procesos de trabajo?
\item ¿Cómo se implementa Kanban en diferentes contextos, como el desarrollo de software, la producción de manufactura, la gestión de proyectos, entre otros?
\item ¿Cuál es el propósito de utilizar tarjetas Kanban en tu proyecto?
\item ¿Cómo vas a estructurar el tablero Kanban en tu proyecto? 
\item ¿Cómo se van a dividir las tareas en las columnas del tablero Kanban? ¿Hay alguna etiqueta o código de colores que utilizarás?
\item ¿Cómo van a ser asignadas las tareas a los miembros del equipo?
\item ¿Cómo se va a hacer seguimiento del progreso de las tareas en el tablero Kanban?
\item ¿Cómo se van a manejar las tareas que no se puedan completar en una iteración? ¿Se van a reprogramar o se van a eliminar?
\item ¿Cómo se van a manejar las tarjetas que necesiten ser reasignadas o modificadas?
\item ¿Cómo se va a informar a los miembros del equipo sobre los cambios en el tablero Kanban? ¿Se utilizará algún sistema de notificaciones?
\item ¿Cómo se va a garantizar la privacidad y seguridad de la información que se maneja en el tablero Kanban?

\end{enumerate}

\end{document}